\documentclass[letterpaper]{article}
\usepackage{aaai24}
\usepackage{times}
\usepackage{helvet}
\usepackage{courier}
\usepackage[hyphens]{url}
\usepackage{graphicx}
\usepackage{amsmath}
\usepackage{booktabs}
\usepackage{multirow}
\usepackage{caption}
\frenchspacing

\pdfinfo{
/TemplateVersion (2024.1)
}

\setcounter{secnumdepth}{0}

\title{ConvStem-ViT: Improving Vision Transformer Patch Embedding for Small-Scale Image Classification}
\author{
    Anonymous Course Project
}
\affiliations{
    Department / University \\
    email@example.com
}

\begin{document}
\maketitle

\begin{abstract}
Vision Transformer (ViT) demonstrates that pure Transformer architectures can
perform competitively on image recognition when trained at scale. However, its
original patch embedding layer converts non-overlapping image patches into
tokens using a direct linear projection, which provides limited local inductive
bias. This project studies a simple modification for small-scale image
classification: replacing the linear patch embedding with a lightweight
convolutional stem while keeping the Transformer encoder unchanged. We evaluate
the method on CIFAR-10 against a vanilla ViT and a ResNet18 baseline under a
shared training protocol. The goal is to test whether early local feature
extraction improves the training behavior and final accuracy of ViT-like models
in a data-limited setting.
\end{abstract}

\section{Introduction}

The Vision Transformer proposed by \citet{dosovitskiy2021image} adapts the
Transformer architecture \citep{vaswani2017attention} to image classification by
dividing an image into fixed-size patches, projecting each patch into a token,
and applying a standard Transformer encoder. This design removes many
hand-crafted image-specific operations and shows strong scalability when large
pre-training datasets are available.

Despite its simplicity, the original ViT patch embedding has weaker local
inductive bias than convolutional neural networks. For small image datasets such
as CIFAR-10 \citep{krizhevsky2009learning}, this may make optimization and
generalization more difficult. Motivated by prior observations that early
convolutions can benefit vision Transformers \citep{xiao2021early}, we explore
a compact ConvStem-ViT variant. The proposed model replaces the direct linear
patch projection with a small stack of strided convolutions, batch
normalization, and GELU activations.

This course project asks a focused question: can a convolutional patch embedding
improve a small ViT under the same Transformer architecture and training
protocol? The contribution is a runnable implementation and a controlled
experimental comparison among ResNet18, vanilla ViT, and ConvStem-ViT on
CIFAR-10.

\section{Related Work}

\paragraph{Vision Transformer.}
ViT treats image patches as visual words and processes them with a Transformer
encoder \citep{dosovitskiy2021image}. Its key innovation is the direct transfer
of the NLP Transformer recipe to image recognition: patch tokenization,
learnable positional embeddings, a class token, multi-head self-attention, and
MLP blocks.

\paragraph{Convolutional networks.}
Residual networks \citep{he2016deep} remain a strong baseline for small-scale
image classification because convolutions naturally encode locality and
translation equivariance. This makes ResNet18 a useful reference point for
evaluating whether the proposed ViT modification narrows the gap between
Transformer and convolutional models.

\paragraph{Hybrid visual architectures.}
Several works study combinations of convolution and attention. In particular,
\citet{xiao2021early} show that early convolutional layers can improve visual
Transformers. Our project follows the same broad motivation but keeps the
implementation intentionally small and course-project friendly.

\section{Method}

\subsection{Original ViT Patch Embedding}

Given an input image $x \in \mathbb{R}^{H \times W \times C}$, ViT splits it
into $N = HW/P^2$ non-overlapping patches, where $P$ is the patch size. Each
flattened patch is projected to a $D$-dimensional token. A learnable class token
and positional embeddings are then added before the sequence is processed by
Transformer encoder layers.

For CIFAR-10, we use $32 \times 32$ images and a patch size of $4$, producing
$8 \times 8 = 64$ patch tokens. This keeps the number of tokens manageable while
preserving the spirit of the original ViT design.

\subsection{ConvStem-ViT}

The proposed ConvStem-ViT replaces the original patch embedding with a
lightweight convolutional stem:

\begin{equation}
z = \mathrm{ConvStem}(x),
\end{equation}

where the stem consists of strided $3 \times 3$ convolutional layers, batch
normalization, and GELU activations. For a patch size of $4$, two stride-2
convolutions reduce the spatial resolution from $32 \times 32$ to $8 \times 8$.
The resulting feature map is flattened into the same number of tokens as the
vanilla ViT. The Transformer encoder depth, embedding dimension, number of
attention heads, and classifier head are unchanged.

This design isolates the effect of patch embedding. Any performance difference
between vanilla ViT and ConvStem-ViT should mainly come from the added local
feature extraction before self-attention.

\section{Experiments}

\subsection{Dataset}

We evaluate on CIFAR-10 \citep{krizhevsky2009learning}, which contains 50,000
training images and 10,000 test images from 10 classes. All images have
resolution $32 \times 32$.

\subsection{Models}

We compare three models:

\begin{itemize}
    \item \textbf{ResNet18}: a convolutional baseline adapted for CIFAR-10.
    \item \textbf{Vanilla ViT}: a small ViT using linear patch projection.
    \item \textbf{ConvStem-ViT}: the same ViT backbone with convolutional patch embedding.
\end{itemize}

\subsection{Training Protocol}

All models are trained using AdamW with learning rate $3 \times 10^{-4}$,
weight decay $0.05$, cosine learning-rate scheduling, batch size 128, and 50
epochs. Data augmentation includes random cropping and random horizontal
flipping. The random seed is fixed to 42.

\begin{table}[t]
\centering
\caption{Training settings used for all experiments.}
\begin{tabular}{lc}
\toprule
Setting & Value \\
\midrule
Dataset & CIFAR-10 \\
Epochs & 50 \\
Batch size & 128 \\
Optimizer & AdamW \\
Learning rate & $3 \times 10^{-4}$ \\
Weight decay & 0.05 \\
Scheduler & Cosine annealing \\
Seed & 42 \\
\bottomrule
\end{tabular}
\label{tab:settings}
\end{table}

\section{Results and Analysis}

Table~\ref{tab:results} should be filled after running the provided experiment
scripts. The expected output files are \texttt{metrics.csv} and
\texttt{summary.json} under each run directory.

\begin{table}[t]
\centering
\caption{Main comparison on CIFAR-10. Replace TBD values after running experiments.}
\begin{tabular}{llcc}
\toprule
Method & Patch Embedding & Accuracy & Params \\
\midrule
ResNet18 & CNN & TBD & TBD \\
Vanilla ViT & Linear & TBD & TBD \\
ConvStem-ViT & Conv stem & TBD & TBD \\
\bottomrule
\end{tabular}
\label{tab:results}
\end{table}

The main hypothesis is that ConvStem-ViT will converge more stably or achieve a
higher test accuracy than vanilla ViT because the convolutional stem captures
local image patterns before global self-attention. If the improvement is small,
the result still provides useful evidence that patch embedding alone may not be
sufficient and that larger data, stronger augmentation, or pre-training may be
needed for ViT-style architectures.

\section{Conclusion}

This project implements and evaluates a simple improvement to the ViT patch
embedding layer. The original ViT demonstrates the power of representing images
as patch tokens and processing them with Transformer encoders. Our ConvStem-ViT
keeps this global attention mechanism but adds local visual feature extraction
at the input stage. The project is designed to be reproducible, lightweight, and
directly connected to the key innovation of the original paper.

\bibliography{references}

\end{document}
